\section{Introduction}
\label{sec:introduction}

Many organisations hold data that they may not pool. Hospitals hold patient
records, banks hold transactions, public agencies hold census and survey data,
and a software service keeps each tenant's data apart from the others.
Analytics and machine learning across such data have to answer the same
questions many times a day:
\begin{itemize}
  \item What may this requester see from this source, and in what form?
  \item May this result, or this model update, leave this site?
  \item May this site take part in this study?
  \item May this new model replace the one in use?
\end{itemize}

Law and professional practice set the rules for these answers. The General
Data Protection Regulation (GDPR) requires purpose limitation and data
minimisation~\citep{eu2016gdpr}. The HIPAA Privacy Rule asks covered entities to
use and disclose only the minimum necessary information~\citep{hhs2026minimum}.
The EU Artificial Intelligence Act requires high-risk systems to keep logs and
to allow human oversight~\citep{eu2024aiact}. Regulators of medical devices ask
manufacturers to plan and control how a learning model may change after it is
cleared~\citep{fda2024pccp}. Each of these rules is checked per request or per
release, not once at design time.

Existing work answers parts of these questions. Statements of responsible-AI
principles and documentation such as model cards describe what a system should
do, but they do not act on a request~\citep{jobin2019global,mitchell2019model}.
Access-control and policy engines decide whether a request is permitted, but
most of them do not decide how much of the data to return or in what
form~\citep{hu2014guide}. Privacy-enhancing technologies provide mechanisms, such
as \(k\)-anonymity and differential privacy, but a person or a separate system
still chooses which mechanism applies to which
request~\citep{sweeney2002kanonymity,dwork2014algorithmic}. Federated learning
frameworks keep raw data at each site and implement the training
protocol~\citep{mcmahan2017communication,kairouz2021advances}, while the
decision about what else may leave a site, and whether the resulting model may
be released, is often left to each project.

GRAILS~\citep{kulkarni2025grails} takes a different route. It separates ethical
restrictions from the functional code of an application. A planner, the
\emph{Ethical-Restriction Rails}, chooses a transformation for each requested
data item from three inputs: the sensitivity of the data, the trust the system
places in the requester, and the granularity of the request (cell, row, column
or table). A runtime component, the \emph{Ethical Guard}, applies that plan to
the data and returns a filtered response. GRAILS compares its planner to the
query planner of a database.

This paper uses GRAILS as its starting point and asks one question of every
approach it surveys: \emph{where do the governance decisions live?} We split
the question into four parts:
\begin{enumerate}
  \item Where is the decision made, and in what form: application code,
        configuration, a policy language, a person, a learned model or a typed
        function?
  \item What does it decide: permit or deny, a transformation, a privacy
        mechanism and its parameters, who may take part, or whether a result
        may be released?
  \item Where is it enforced: at a central store, at each data source, at the
        point where data leaves, or at model release?
  \item What evidence does it leave, and can the decision be replayed later?
\end{enumerate}

We then describe the design idea behind Dagents, an open-source framework that
extends the GRAILS split from one request boundary in one application to every
boundary of a study that runs across several data sources. Dagents runs a local
agent at each source and a global agent that coordinates them. It makes every
governance and federation decision with pure functions written in OCaml, over
algebraic data types, and enforces each decision at four boundaries: before a
site accepts a training job, before it reads its records, before anything
leaves it, and before the coordinator releases a model.

This paper makes four contributions:
\begin{enumerate}
  \item A survey of \SurveyWorks{} works and industry systems in seven families,
        each described by the four questions above (\cref{sec:survey}).
  \item A synthesis that shows which parts of the lifecycle of a multi-source
        study each family governs, and which parts no family governs as a whole
        (\cref{sec:synthesis}).
  \item A statement of the Dagents design as ten principles, each tied to code
        that implements it (\cref{sec:design}).
  \item A comparison of that design with GRAILS and with each surveyed family,
        a case study on 1,014 real patient records from three hospitals in the
        International Stroke Trial, and the limits of the design
        (\cref{sec:comparison,sec:case-study,sec:limitations}).
\end{enumerate}

We built Dagents, so we are not neutral judges of it. To reduce this bias, the
survey records, for every family, what it does that Dagents does not, and
\cref{sec:limitations} lists the claims this paper does not make.
